\documentclass[11pt]{article}
\usepackage[margin=1in]{geometry}
\usepackage{amsmath,amssymb,amsthm,enumitem,hyperref}
\hypersetup{colorlinks=true,linkcolor=blue,urlcolor=blue,citecolor=blue}
\newtheorem{theorem}{Theorem}
\newcommand{\E}{\mathbb E}
\newcommand{\Q}{\mathcal Q}
\newcommand{\des}{\operatorname{des}}
\title{The descent formula for preorder polytopes\\\large MathDB 374958}
\author{Principia Math In Progress}
\date{September 18, 2026}
\begin{document}
\maketitle
\begin{abstract}
For every preorder on an arbitrarily totally ordered $n$-element set,
we prove that the $h^*$-polynomial of its preorder polytope is the
reverse-descent enumerator of words meeting every order ideal at
least as often as its cardinality. This proves Conjecture 4.4 of
Athanasiadis and Chapoton. We use bivariate Ehrhart reciprocity,
their established counting formula, and an explicit multiset argument.
\end{abstract}

\section*{Task 1. Full solution}
Let $E$ be a finite set with $|E|=n\ge1$, equipped with a preorder
$\tau$ and, separately, an arbitrary total order. There is no
compatibility assumption between these two orders. Set
\[
 \Q_\tau=\left\{x\in\mathbb R_{\ge0}^E:
       \sum_{e\in I}x_e\le |I|\text{ for every order ideal }I\right\}.
\]
Write $\tau^*$ for the reversed preorder, and
$P_{\tau^*}=\Q_{\tau^*}\cap\mathbb Z^E$.
Let $W_\tau$ consist of the words $w\in E^n$ such that the letters
from each order ideal $I$ occur at least $|I|$ times in $w$.
The number $\des(w)$ counts indices $1\le j<n$ with $w_j>w_{j+1}$
in the chosen total order.

\begin{theorem}\label{thm:main}
For every such preorder and total order,
\[
 \boxed{h^*(\Q_\tau,t)=\sum_{w\in W_\tau}t^{n-1-\des(w)}.}
\]
\end{theorem}

\subsection*{A counting identity from reciprocity}
We give the needed reciprocity argument in full. Let
\[
 \Q_\tau(r,s)=\left\{x\in\mathbb R_{\ge0}^E:
       \sum_{e\in I}x_e\le r|I|+s
       \text{ for every nonempty order ideal }I\right\}.
\]
By \cite[Propositions 3.2--3.3 and 7.1]{AC}, for nonnegative
integers $r,s$ this is a lattice polytope, also expressible as an
integral Minkowski sum of simplices. Its lattice-point count is the
polynomial $\E_\tau(r,s)$, of total degree $n$.

For positive integers $r,s$, the polytope is full dimensional
(it contains $[0,r]^E$). Its interior integer points satisfy
$x_e\ge1$ and $\sum_{e\in I}x_e\le r|I|+s-1$ for every
nonempty order ideal. Subtracting the all-ones vector therefore
bijects these points with $\Q_\tau(r-1,s-1)\cap\mathbb Z^E$.
All inequality normals are nonzero; redundant inequalities have
no effect on this strict-interior characterization.

For $P=\Q_\tau(r,s)$ one has $kP=\Q_\tau(kr,ks)$ for
nonnegative integers $k$, so its Ehrhart polynomial equals
$\E_\tau(kr,ks)$ as a polynomial in $k$. Ordinary Ehrhart
reciprocity at $k=-1$ and the translation just described imply
\begin{equation}\label{eq:rec}
 \E_\tau(-u,-v)=(-1)^n\E_\tau(u-1,v-1).
\end{equation}
Initially this holds for positive integer pairs $(u,v)$; it is a
polynomial identity because a polynomial vanishing on that entire
grid is zero (first fix $v$ and vary $u$, then vary $v$).

For $a=(a_e)_{e\in E}$ write $|a|=\sum_e a_e$ and $a_0=n-|a|$.
The established formula \cite[Lemma 4.1 and Equation (25)]{AC} is
\begin{equation}\label{eq:source}
 \E_\tau(u,v)=\sum_{a\in P_{\tau^*}}
       \binom{v+a_0}{a_0}\prod_{e\in E}\binom{u+a_e-1}{a_e}.
\end{equation}
Here $a_0\ge0$, and binomial coefficients with fixed nonnegative
lower argument are interpreted as polynomials in the upper argument.
Putting $(u,v)=(-m-1,-1)$ in \eqref{eq:source}, every term with
$a_0>0$ vanishes, since $\binom{a_0-1}{a_0}=0$.
For the remaining terms $|a|=n$, use
\[
 \binom{a_e-m-2}{a_e}=(-1)^{a_e}\binom{m+1}{a_e}.
\]
Together with \eqref{eq:rec}, this gives, for all $m\ge0$,
\begin{equation}\label{eq:count}
 \operatorname{Ehr}(\Q_\tau,m)=\E_\tau(m,0)
 =\sum_{\substack{a\in P_{\tau^*}\\|a|=n}}
        \prod_{e\in E}\binom{m+1}{a_e}.
\end{equation}
This implication of bivariate reciprocity was already noted,
conditionally, in \cite[Remark 7.3]{AC}; equation \eqref{eq:rec}
supplies the missing premise.

\subsection*{An elementary multiset identity}
For a nonnegative integer vector $a$ with $|a|=n$, let $S_a$
be all words having $a_e$ copies of each $e$. We prove
\begin{equation}\label{eq:word}
 \sum_{m\ge0}\left(\prod_e\binom{m+1}{a_e}\right)t^m
 =\frac{\displaystyle\sum_{w\in S_a}t^{n-1-\des(w)}}{(1-t)^{n+1}}.
\end{equation}
Choose, for each letter $e$, a subset of $a_e$ bins from the
$m+1$ ordered bins $0,1,\ldots,m$. Place one copy of $e$ in each
chosen bin. The number of choices is the product on the left.
Within each bin, read its distinct letters in strictly decreasing
order, and concatenate the bins. This produces a word $w\in S_a$
and a weakly increasing sequence of bin indices
$0\le b_1\le\cdots\le b_n\le m$.

At each index $j$ for which $w_j\le w_{j+1}$, the bin index must
strictly increase. Conversely these conditions suffice: a constant
bin segment is strictly decreasing and contains no repeated letter.
Thus this construction is a bijection. The number of required
strict increases is $c(w)=n-1-\des(w)$. Subtract from $b_j$ the
number of required increases before position $j$. When $m\ge c(w)$
this gives a weakly increasing sequence of length $n$ in
$\{0,\ldots,m-c(w)\}$, counted by
$\binom{m-c(w)+n}{n}$. There are no such sequences for $m<c(w)$.
Consequently the contribution of $w$ to the generating function is
\[
 \sum_{m\ge c(w)}\binom{m-c(w)+n}{n}t^m
 =\frac{t^{c(w)}}{(1-t)^{n+1}},
\]
which proves \eqref{eq:word}, including repeated letters and $m=0$.

\begin{proof}[Proof of Theorem~\ref{thm:main}]
A vector $a$ of total sum $n$ belongs to $P_{\tau^*}$ exactly
when $\sum_{e\in J}a_e\le |J|$ for every order filter $J$ of
$\tau$. Taking complements $I=E\setminus J$, these conditions
are equivalent to $\sum_{e\in I}a_e\ge |I|$ for every order
ideal $I$ of $\tau$. Therefore the sets $S_a$ indexed in
\eqref{eq:count} partition $W_\tau$.
Sum \eqref{eq:word} over those vectors and use \eqref{eq:count}.
The definition
\[
 \sum_{m\ge0}\operatorname{Ehr}(\Q_\tau,m)t^m
 =\frac{h^*(\Q_\tau,t)}{(1-t)^{n+1}}
\]
now yields the desired identity.
\end{proof}

\subsection*{Verification and attribution}
Athanasiadis and Chapoton supplied the counting formula and
proposed the implication from reciprocity. The complete
reciprocity proof is reproduced above from our companion solution
to MathDB 374965. The elementary multiset argument makes the
descent conversion explicit. The earlier result \cite{AXY} treats
a special family of arbors. The later preorder result in
\cite[Corollary 5.4]{Y} concerns the reflexive polytope
$\Q_\tau+\Delta_E-\mathbf1$, whose word set includes an extra
letter; it is a different statement.

All 26 pages of the current primary were read. The accompanying
verifier compares direct lattice counts with word enumeration for
all 389 labeled preorders of sizes 1--4, checks 50,069 words for
the multiset identity, and tests all 120 total orderings in the
source's five-element example. Bounded checks found no prior full
target claim or identified journal final. Review is by the
authoring agent only, without independent peer review or a proof assistant.

\section*{Tasks 2--5. Proposed follow-up directions}
\begin{enumerate}[label=\textbf{Task \arabic*.},start=2,leftmargin=*]
\item Construct an explicit bijection directly from lattice points
of $m\Q_\tau$ to the bin assignments in \eqref{eq:count}.
\item Find a geometric realization of descent-set refinements of
the formula, retaining the positions of descents rather than their number.
\item Derive colored-word analogues for integral weights assigned
to the principal-filter simplex summands.
\item Characterize extremal $h^*$-coefficients among preorders of
fixed size using the ideal constraints on the word set.
\end{enumerate}
\begin{thebibliography}{9}
\raggedright
\bibitem{AC} C. A. Athanasiadis and F. Chapoton,
\emph{Polytopes and posets associated to preorders},
\href{https://arxiv.org/abs/2605.26916v1}{arXiv:2605.26916v1}, 2026.
\bibitem{AXY} C. A. Athanasiadis, Q. Xiao, and X. Yan,
\emph{Lattice point enumeration of some arbor polytopes},
\href{https://arxiv.org/abs/2603.11654v1}{arXiv:2603.11654v1}, 2026.
\bibitem{Y} X. Yan, \emph{Variations of colored multiset Eulerian
polynomials and applications},
\href{https://arxiv.org/abs/2608.15682v2}{arXiv:2608.15682v2}, 2026.
\end{thebibliography}
\end{document}
